\chapter*{Введение}
\addcontentsline{toc}{chapter}{Введение}

В ходе расчетно-графической работы была спроектирована и реализована база данных для управления SIM-картами и телефонными номерами мобильного оператора. Такая система нужна для учета абонентов, тарифов, выданных номеров, SIM-карт, блокировок и истории операций.

Цель работы состояла в том, чтобы разработать связанную реляционную структуру базы данных, реализовать ее в PostgreSQL, наполнить тестовыми данными и подготовить набор SQL-запросов для проверки работы системы. Дополнительно была выполнена контейнеризация базы данных с помощью Docker Compose, чтобы проект можно было быстро запустить на другом компьютере.

В работе были выполнены следующие задачи:

\begin{enumerate}
    \item Определена предметная область и основные процессы системы.
    \item Спроектирована структура базы данных из связанных таблиц.
    \item Реализован файл \texttt{schema.sql} для создания таблиц, ключей, ограничений и индексов.
    \item Подготовлен файл \texttt{data.sql} с тестовыми данными.
    \item Составлен файл \texttt{queries.sql} с десятью функциональными запросами.
    \item Построена ER-диаграмма базы данных.
    \item Настроен запуск PostgreSQL через \texttt{docker-compose.yaml}.
    \item Проверена работоспособность запросов в контейнере PostgreSQL.
\end{enumerate}

Отчет составлен по проделанной работе. В нем описана логика выбранных таблиц, связи между ними, ограничения целостности, запуск проекта и результаты выполнения запросов.
